% !TeX root = ../nc-quadric-moduli.tex
% 2026-11-22 / v0.1（表示商の特殊fiber証明を論文Bへ移管）
\subsection{良モジュライ射の平坦性と特殊ファイバー}
\label{subsec:moduli-noncubic-fibers}

本節は表示商 $\mathfrak N^{\mathrm{pres}}=[P/G]$ そのものの計算であり、
その内在的被約境界との相対同値は分類原稿
\cite{QuadricIntrinsic2026} の結果を別途入力とする。
$Q^{\mathrm{nd}}=\operatorname{Sym}^2 V\setminus V(\operatorname{disc})$、
$P=Q^{\mathrm{nd}}\times Q^{\mathrm{nd}}$ とする。
\begin{proposition}[平坦な良モジュライ射と特殊 fiber]
\label{prop:moduli-noncubic-good-moduli-fibers}
$\chi$ は平坦かつ有限表示である。$\xi\in k$ に対し
$\mathfrak N^{\mathrm{pres}}_\xi=\chi^{-1}(\xi)$ を scheme-theoretic fiber とすると、次が成り立つ。
\begin{enumerate}
\item $\xi\ne0,1$ の fiber は smooth である。
\item $\xi=0$ の fiber は smooth な被約 Cartier divisor の二重構造である。
\item $\xi=1$ の fiber は reduced である。その幾何点は、一根共有の類と二根対一致の類の二つである。
一根共有の類は smooth な開部分をなし、その閉包は一致類を含む。
一致類の近傍で、smooth chart 上の fiber 方程式は $ac=0$ となる。
\end{enumerate}
従って $\chi$ が smooth でない点は、$\xi=0$ の全点と $\xi=1$ の一致類に限られる。
\end{proposition}

\begin{proof}
まず固定した第一形式への帰着を smooth な被覆として実現する。
$\operatorname{GL}(V)$ は $Q^{\mathrm{nd}}$ に推移的に作用し、
$uv$ の安定化群は $O(uv)\simeq\mathbb G_m\rtimes C_2$ で smooth である。
従って軌道射 $\operatorname{GL}(V)\to Q^{\mathrm{nd}}$ は smooth 全射である。
例えば単位元における微分は、二つの非対角成分から $u^2,v^2$ を、
対角成分の和から $uv$ を得るため全射であり、群作用により全点で同様である。
その基底変更に同型な射
\[
 \operatorname{GL}(V)\times Q^{\mathrm{nd}}\longrightarrow P,
 \qquad (A,g)\longmapsto
 \bigl(\operatorname{Sym}^2(A)(uv),\operatorname{Sym}^2(A)g\bigr)
\]
も smooth 全射である。
さらに $P\to[P/G]$ は smooth 全射なので、これらの合成上で fiber の性質を調べられる。

$g=au^2+buv+cv^2$、$\Delta=b^2-4ac$ と書けば、この被覆上で
\[
 \chi=\frac{b^2}{\Delta},\qquad \Delta\in\OO^\times.
\]
従って fiber の方程式は
\begin{equation}\label{eq:moduli-noncubic-fiber-quadric}
 (1-\xi)b^2+4\xi ac=0,\qquad \Delta\ne0
\end{equation}
である。$\xi\ne0,1$ なら、この二次形式は三変数の非退化二次形式であり、
その唯一の affine 特異点である原点は $\Delta\ne0$ により除かれる。
従って fiber は smooth である。

$\xi=0$ では式は $b^2=0$ となる。
被約部分 $b=0$ 上では $ac$ が可逆なので、これは smooth な Cartier divisor の二重構造である。
この divisor は $P$ 上では判別式の極化 $D(f_1,f_2)=0$ で与えられる。
$D$ は $G$ の半不変式であり、分母 $\Delta(f_1)\Delta(f_2)$ は可逆であるため、
商スタック上でも $\operatorname{div}(\chi)=2\mathfrak E_0$ と書ける。
ここで $\mathfrak E_0$ は smooth な被約 Cartier divisor である。

$\xi=1$ では $ac=0$ となり、$\Delta\ne0$ は $b$ が可逆であることと同値である。
これは二つの座標軸の通常交差と可逆な $b$ の積なので reduced である。
$(a,c)\ne(0,0)$ なら二形式は一根だけを共有し、fiber は smooth である。
$a=c=0$ は二形式が比例する場合であり、fiber はそこで交差特異点を持つ。
$b$ を一に正規化した横断座標では
\begin{equation}\label{eq:moduli-noncubic-chi-coincident-node}
 \chi-1=\frac{4ac}{1-4ac}.
\end{equation}
$a$ 軸または $c$ 軸の原点への特殊化が、一根共有の類から一致類への特殊化である。
二つの軸は $u,v$ の交換で同じ軌道へ移るので、幾何点の類は記述した二つで尽くされる。

平坦性は $P$ 上で確認できる。
$P$ は整 affine スキームであり、$\chi$ は非定数なので $k[\chi]\to k[P]$ は単射である。
従って $k[P]$ は PID $k[\chi]$ 上 torsion-free、従って flat である。
有限表示性は有限型 $k$-スキーム・スタックの間の射であることから従う。
平坦性は smooth 全射 $P\to\mathfrak N^{\mathrm{pres}}$ から降下する。
最後の smooth locus の記述は、平坦有限表示性と上で判定した幾何 fiber の smooth locus から従う。
\end{proof}

\begin{remark}[resultant による共有根の方程式]
二つの二元二次形式に対し、標準的な Sylvester resultant を用いると
\[
 D(f_1,f_2)^2-\Delta(f_1)\Delta(f_2)=4\operatorname{Res}(f_1,f_2).
\]
従って $\chi=1$ は共有根の resultant divisor である。
二つの形式自身は非退化であるが、共有根が二つになる一致類では、この divisor の局所方程式が $ac=0$ となる。
\end{remark}